\chapter{Volba technologií}\label{chap:technologies}

\begin{chapterabstract}
	Tato kapitola srovnává technologie a~frameworky vhodné pro implementaci jednotlivých částí aplikace. Pro každou vrstvu (backend, frontend, mobilní aplikace) jsou porovnány dostupné alternativy s~vyhodnocením výhod a~nevýhod. Kapitola je zakončena přehledem dalších klíčových knihoven použitých v~projektu.
\end{chapterabstract}

Vyvíjený systém se skládá ze tří hlavních komponent: serverové části (backend), webové aplikace (frontend) a~mobilní aplikace pro inventuru (Android). Každá z~těchto komponent vyžaduje volbu vhodné technologie, přičemž je nutné zohlednit vzájemnou kompatibilitu a~dostupnost knihoven. Všechny zvažované technologie jsou open-source a~bezplatné. Při volbě byla zohledněna i~předchozí zkušenost autora, která umožnila efektivní vývoj.

\section{Backend}

Backend aplikace zajišťuje veškerou business logiku, komunikaci s~databází, autentizaci uživatelů a~generování exportů. Jedná se o~REST API, které obsluhuje jak webovou aplikaci, tak mobilní aplikaci. Pro jeho implementaci byly zvažovány tři technologie.

\subsection{Spring Boot (Java)}

Spring Boot je framework postavený nad platformou Spring, který podle oficiální dokumentace \uv{pomáhá vytvářet samostatné, produkčně kvalitní aplikace založené na Springu s~minimálním úsilím} (překlad vlastní)~\cite{springboot}. Framework zastává tzv. opinionated přístup -- poskytuje rozumné výchozí konfigurace, které lze v~případě potřeby přepsat.

Mezi hlavní přednosti Spring Boot patří rozsáhlý ekosystém. Spring Data JPA~\cite{springdatajpa} nabízí abstrakci nad databázovou vrstvou pomocí Jakarta Persistence API, Spring Security~\cite{springsecurity} poskytuje komplexní řešení autentizace a~autorizace včetně podpory JWT tokenů a~OAuth2. Framework dále disponuje vestavěným webovým serverem (Tomcat), automatickou konfigurací na základě přítomných závislostí a~robustní podporou pro testování prostřednictvím JUnit~5 a~Mockito.

Nevýhodou Spring Boot je vyšší paměťová náročnost typická pro aplikace běžící na JVM a~delší doba studeného startu oproti interpretovaným jazykům. Pro aplikace s~nízkou frekvencí změn v~nasazení, jako je tento projekt, však tato nevýhoda není podstatná.

\subsection{NestJS (TypeScript)}

NestJS~\cite{nestjs} je Node.js framework nabízející modulární architekturu inspirovanou Angularem, vestavěný systém Dependency Injection a~podporu TypeScriptu jako prvotřídního jazyka. Je postaven nad minimalistickým Express.js~\cite{expressjs}, ale na rozdíl od něj poskytuje kompletní strukturu aplikace. Pro práci s~databází lze využít TypeORM nebo Prisma.

Hlavní výhodou je možnost sdílení TypeScriptových typů mezi backendem a~frontendem. Pro tento projekt je nevýhodou absence knihoven srovnatelných s~iText (generování PDF) a~EasyExcel (export do XLSX) v~ekosystému Node.js -- dostupné alternativy (např.\ PDFKit, ExcelJS) nenabízejí stejnou úroveň podpory pro složitější tabulkové reporty a~lokalizaci.

\subsection{Django (Python)}

Django~\cite{django} je plnohodnotný webový framework pro Python, zastávající filozofii \uv{batteries included} -- součástí je vestavěný ORM, administrační rozhraní, systém autentizace i~ochrana proti běžným bezpečnostním hrozbám (CSRF, SQL injection).

Automaticky generované administrační rozhraní je výhodou pro projekty vyžadující rychlý prototyping. Nevýhodou Djanga je dynamický typový systém Pythonu, který oproti Javě nenabízí statickou kontrolu typů v~době kompilace, což může vést k~obtížněji odhalitelným chybám. Ekosystém knihoven pro generování PDF a~ZPL etiket je v~Pythonu rovněž méně rozvinutý než v~prostředí JVM.

\subsection{Srovnání a~zdůvodnění volby}

Tabulka~\ref{tab:backend-comparison} shrnuje srovnání analyzovaných backend technologií.

\begin{table}[htbp]\centering\footnotesize
\begin{tabular}{l|c|c|c}
	\textbf{Kritérium} & \textbf{Spring Boot} & \textbf{NestJS} & \textbf{Django} \tabularnewline \hline
	Jazyk & Java & TypeScript & Python \tabularnewline
	Typová bezpečnost & silná & silná & dynamická \tabularnewline
	ORM & JPA (Hibernate) & TypeORM / Prisma & vlastní ORM \tabularnewline
	Bezpečnostní framework & Spring Security & Passport.js & vestavěný \tabularnewline
	Podpora PDF/Excel/ZPL & iText, EasyExcel & omezená & omezená \tabularnewline
\end{tabular}
\caption{Srovnání backend technologií}\label{tab:backend-comparison}
\end{table}

Pro implementaci backendu byl zvolen \textbf{Spring Boot} z~následujících důvodů:
\begin{itemize}
	\item Robustní ekosystém pro enterprise aplikace -- Spring Data JPA pro databázovou vrstvu, Spring Security pro autentizaci a~autorizaci.
	\item Dostupnost kvalitních knihoven pro generování PDF (iText), Excelu (EasyExcel) a~ZPL etiket -- klíčová funkčnost pro tento projekt.
	\item Silná typová bezpečnost jazyka Java, která usnadňuje údržbu a~refaktoring.
	\item Konzistence s~mobilní aplikací, která je rovněž implementována v~Javě.
\end{itemize}


\section{Frontend}

Webová aplikace slouží jako hlavní rozhraní pro správu majetkových karet, číselníků, uživatelů a~vyhodnocení inventurních dat. Jedná se o~jednostránkovou aplikaci (SPA) komunikující s~backendem prostřednictvím REST API. Pro její implementaci byly zvažovány tři přední JavaScriptové technologie.

\subsection{React}

React je podle své dokumentace \uv{knihovna pro tvorbu webových a~nativních uživatelských rozhraní} (překlad vlastní)~\cite{react}. Na rozdíl od Angularu se nejedná o~kompletní framework, ale o~knihovnu zaměřenou výhradně na vrstvu uživatelského rozhraní. Výběr ostatních nástrojů (routing, state management, formuláře) je ponechán na vývojáři.

React využívá komponentovou architekturu a~deklarativní přístup prostřednictvím JSX syntaxe. Systém Hooks (\texttt{useState}, \texttt{useEffect}) umožňuje správu stavu a~vedlejších efektů ve funkcionálních komponentách. Ekosystém Reactu je největší ze srovnávaných technologií -- zahrnuje knihovny jako MUI~\cite{mui} pro komponentové UI, TanStack Query~\cite{tanstackquery} pro správu serverového stavu, React Hook Form~\cite{reacthookform} pro formuláře či Zustand~\cite{zustand} pro globální stav.

Nevýhodou Reactu je nutnost skládat aplikaci z~mnoha nezávislých knihoven, což vyžaduje zkušenost s~výběrem a~integrací těchto nástrojů.

\subsection{Angular}

Angular~\cite{angular} je plnohodnotný, opinionated framework vyvíjený týmem ve společnosti Google. Angular poskytuje kompletní sadu nástrojů přímo v~základním balíčku -- vlastní router, HTTP klient, systém formulářů, Dependency Injection a~build toolchain.

TypeScript je v~Angularu povinný, což zajišťuje silnou typovou bezpečnost. Framework dále nabízí reaktivní systém Angular Signals a~vestavěnou podporu pro Server-Side Rendering. Nevýhodou je strmější křivka učení a~větší režie pro menší projekty, kde plnohodnotný framework může být zbytečně komplexní.

\subsection{Vue.js}

Vue.js~\cite{vuejs} je progresivní framework navržený s~důrazem na inkrementální adoptovatelnost -- lze jej použít jak jako jednoduchou knihovnu pro obohacení existujících stránek, tak jako základ pro plnohodnotnou SPA.

Vue nabízí dvě API rozhraní: Options API (přívětivější pro začátečníky) a~Composition API (doporučené pro větší aplikace). Ekosystém obsahuje oficiální řešení pro routing (Vue Router) a~state management (Pinia). Oproti Reactu je ekosystém UI knihoven menší, i~když dostupné jsou například Vuetify nebo Quasar.

\subsection{Srovnání a~zdůvodnění volby}

Tabulka~\ref{tab:frontend-comparison} shrnuje srovnání analyzovaných frontend technologií.

\begin{table}[htbp]\centering\footnotesize
\begin{tabular}{l|c|c|c}
	\textbf{Kritérium} & \textbf{React} & \textbf{Angular} & \textbf{Vue.js} \tabularnewline \hline
	Typ & knihovna & framework & framework \tabularnewline
	TypeScript & volitelný & povinný & volitelný \tabularnewline
	State management & Zustand, Redux & Signals, NgRx & Pinia \tabularnewline
	UI knihovny & MUI, shadcn & Angular Material & Vuetify, Quasar \tabularnewline
	Přístup & neopinionated & opinionated & progresivní \tabularnewline
\end{tabular}
\caption{Srovnání frontend technologií}\label{tab:frontend-comparison}
\end{table}

Pro implementaci frontendu byla zvolena knihovna \textbf{React} společně s~jazykem
\textbf{TypeScript}~\cite{typescript} z~následujících důvodů:
\begin{itemize}
	\item Největší ekosystém UI knihoven -- MUI~\cite{mui} poskytuje hotové komponenty odpovídající Material Design, včetně tabulek, dialogů a~formulářových prvků potřebných pro správu dat.
	\item Flexibilita při výběru knihoven -- možnost zvolit optimální nástroj pro každou oblast (TanStack Query pro data fetching, React Hook Form + Zod pro formuláře, Zustand pro globální stav).
\end{itemize}

Jako build nástroj byl zvolen \textbf{Vite}~\cite{vite}, který nabízí rychlý vývojový server s~Hot Module Replacement a~optimalizovaný produkční build.


\section{Mobilní aplikace}

Mobilní aplikace slouží k~provádění inventury v~terénu -- pracovník skenuje Data Matrix kódy na majetku pomocí čtečky Zebra TC22 a~zaznamenává skutečné umístění položek. Aplikace pracuje offline a~data synchronizuje se serverem, jakmile je dostupné připojení k~síti. Cílovým zařízením je výhradně Android.

\subsection{Nativní Android (Java)}

Nativní vývoj pro Android~\cite{android} v~jazyce Java poskytuje přímý přístup ke všem API platformy bez jakékoliv abstrakční vrstvy. Pro lokální ukládání dat je k~dispozici knihovna Room~\cite{androidroom}, která podle oficiální dokumentace \uv{poskytuje abstrakční vrstvu nad SQLite umožňující plynulý přístup k~databázi při využití plného potenciálu SQLite} (překlad vlastní). Room ověřuje SQL dotazy v~době kompilace a~minimalizuje boilerplate kód prostřednictvím anotací.

Pro HTTP komunikaci lze využít knihovnu OkHttp~\cite{okhttp}, která podporuje synchronní i~asynchronní požadavky. Hlavní výhodou nativního přístupu je plná kontrola nad lifecycle komponent a~přímá kompatibilita se vstupem z~hardwarového skeneru, který posílá data jako klávesnicové události.

Nevýhodou je omezení na jedinou platformu a~větší množství boilerplate kódu oproti modernějším alternativám.

\subsection{Kotlin}

Kotlin~\cite{kotlin} je oficiálně doporučený jazyk pro vývoj Android aplikací. Nabízí modernější syntaxi než Java -- \texttt{null} safety na úrovni typového systému, korutiny pro asynchronní programování a~méně boilerplate kódu. Kotlin je plně interoperabilní s~Javou, což umožňuje postupnou migraci existujícího kódu.

Pro tento projekt by Kotlin nepřinesl zásadní výhody. Aplikace je architektonicky jednoduchá (jedna Activity, tři fragmenty) a~nevyužívá složité asynchronní vzory, pro které jsou korutiny navrženy. Zavedení dalšího jazyka do projektu by zvýšilo komplexitu bez odpovídajícího přínosu.

\subsection{Flutter / React Native}

Flutter~\cite{flutter} je multiplatformní framework od Google, který používá jazyk Dart a~vlastní vykreslovací engine obcházející nativní UI komponenty platformy. React Native~\cite{reactnative} vykresluje JavaScriptový kód prostřednictvím nativních komponent; doporučeným meta-frameworkem je Expo.

Obě technologie nabízejí multiplatformní vývoj, který však v~tomto projektu není potřeba -- cílovým zařízením je výhradně Zebra TC22 se systémem Android. Zásadním problémem je omezená podpora specifického hardwaru: vstup ze skeneru přichází jako události hardwarové klávesnice zakončené klávesou Enter nebo Tab, což vyžaduje přímý přístup k~nativním listenerům (\texttt{OnKeyListener}, \texttt{OnEditorActionListener}). Abstrakční vrstva multiplatformních frameworků by tento přístup komplikovala.

\subsection{Srovnání a~zdůvodnění volby}

Tabulka~\ref{tab:mobile-comparison} shrnuje srovnání analyzovaných mobilních technologií.

\begin{table}[htbp]\centering\footnotesize
\begin{tabular}{l|c|c|c}
	\textbf{Kritérium} & \textbf{Java (nativní)} & \textbf{Kotlin} & \textbf{Flutter / RN} \tabularnewline \hline
	Platformy & Android & Android (+ KMP) & Android + iOS \tabularnewline
	Přístup k~HW skeneru & přímý & přímý & přes bridge \tabularnewline
	Offline DB & Room (SQLite) & Room (SQLite) & SQLite plugin \tabularnewline
	Složitost & nízká & nízká & střední \tabularnewline
\end{tabular}
\caption{Srovnání mobilních technologií}\label{tab:mobile-comparison}
\end{table}

Pro implementaci mobilní aplikace byl zvolen \textbf{nativní Android v~Javě} z~následujících důvodů:
\begin{itemize}
	\item Přímý přístup k~hardwarovému skeneru Zebra -- vstup přichází jako klávesnicové události, které je nutné zpracovat prostřednictvím nativních listenerů.
	\item Plná kontrola nad lifecycle komponent -- aplikace pracuje offline s~lokální Room databází a~vyžaduje přesnou správu synchronizace dat.
	\item Jednoduchost -- pro aplikaci s~jednou Activity a~třemi fragmenty není potřeba složitý framework.
	\item Konzistence s~backendem -- oba projekty sdílejí jazyk Java, což usnadňuje přechod mezi jednotlivými částmi projektu.
\end{itemize}


\section{Databázový systém}

Systém využívá relační databázi pro ukládání majetkových karet, číselníků, uživatelů a~inventurních záznamů. Na mobilním zařízení je navíc potřeba lokální databáze pro offline režim. Pro produkční nasazení byly zvažovány dva databázové systémy.

\subsection{PostgreSQL}

PostgreSQL~\cite{postgresql} je objektově-relační databázový systém s~více než 35 lety aktivního vývoje. Nabízí podporu komplexních datových typů (JSONB, pole), fulltextové vyhledávání, rozšiřitelnost prostřednictvím vlastních funkcí (např.\ \texttt{unaccent} pro vyhledávání bez diakritiky) a~vysokou spolehlivost.

\subsection{MySQL / MariaDB}

MySQL~\cite{mysql} je nejrozšířenější open-source relační databáze s~širokou podporou u~hostingových poskytovatelů. MariaDB~\cite{mariadb} je komunitní fork MySQL se zpětnou kompatibilitou, který MySQL v~open-source prostředí postupně nahrazuje. Pro účely tohoto projektu by oba systémy poskytovaly dostatečnou funkčnost. PostgreSQL byl zvolen díky podpoře funkce \texttt{unaccent} pro vyhledávání bez diakritiky a~předchozí zkušenosti autora.

\subsection{H2 a~Room}

Pro automatizované testování backendu byla zvolena in-memory databáze \textbf{H2}~\cite{h2database} v~PostgreSQL compatibility mode. H2 je kompatibilní s~JPA a~nevyžaduje žádnou externí infrastrukturu, testy tak běží izolovaně a~reproducibilně. Testy ověřují business logiku a~JPA dotazy; SQL specifika PostgreSQL (např.\ funkce \texttt{unaccent}) jsou testována prostřednictvím integračních testů.

Na mobilním zařízení je využita knihovna \textbf{Room}~\cite{androidroom} jako abstrakce nad SQLite. Room poskytuje ověřování SQL dotazů v~době kompilace, automatické mapování mezi Java objekty a~databázovými tabulkami prostřednictvím anotací (\texttt{@Entity}, \texttt{@Dao}, \texttt{@Query}) a~podporu migrací.

\subsection{Zdůvodnění volby}

Pro produkční nasazení byl zvolen \textbf{PostgreSQL} díky zralosti, spolehlivosti a~výborné integraci se Spring Data JPA. Všechny zvažované databázové systémy jsou open-source a~bezplatné, takže náklady na licence nebyly rozhodovacím faktorem. Kombinace PostgreSQL (produkce), H2 (testy) a~Room/SQLite (mobilní zařízení) pokrývá potřeby všech tří prostředí.


\section{Další klíčové knihovny}

Kromě hlavních technologií popsaných výše projekt využívá řadu specializovaných knihoven, které výrazně zefektivňují vývoj.

\subsection{Backend knihovny}

\textbf{MapStruct}~\cite{mapstruct} je generátor kódu pro mapování mezi Java beany (typicky DTO a~entity). Na rozdíl od reflexe-based řešení generuje mapovací kód v~době kompilace, což zajišťuje typovou bezpečnost a~vyšší výkon.

\textbf{Project Lombok}~\cite{lombok} redukuje boilerplate kód prostřednictvím anotací~-- \texttt{@Getter}, \texttt{@Setter} a~\texttt{@RequiredArgsConstructor} generují gettery, settery a~konstruktory automaticky. Lombok je nakonfigurován tak, aby generovaný kód byl označen anotací \texttt{@Generated}, čímž je vyloučen z~měření pokrytí kódu testy.

\subsection{Frontend knihovny}

\textbf{TanStack Query}~\cite{tanstackquery} zajišťuje správu serverového stavu -- automatické cachování, invalidaci a~refetching dat z~API. Všechna data ve webové aplikaci jsou načítána výhradně prostřednictvím TanStack Query hooků.

\textbf{React Hook Form}~\cite{reacthookform} v~kombinaci s~validační knihovnou \textbf{Zod}~\cite{zod} zajišťuje správu formulářů s~typově bezpečnou validací. Validační schémata definovaná v~Zod slouží zároveň jako dokumentace očekávaného formátu dat.

\textbf{Zustand}~\cite{zustand} je minimalistická knihovna pro globální state management, využitá pro správu autentizačního stavu (uživatel, JWT token, expirace).

\textbf{MUI}~\cite{mui} (Material UI) poskytuje hotové React komponenty implementující Material Design -- tabulky, dialogová okna, formulářové prvky a~navigaci. Díky MUI bylo možné rychle vytvořit konzistentní a~přístupné uživatelské rozhraní.


